\begin{lemma}[Lemma 4.2]
Let $n\ge2$, $p>1$, $0<\tau\le n$, $N\ge0$ and $d>0$, and suppose the
Sobolev--Poincar\'e inequality on balls holds in dimension $n$
(\noderef{CitedSobolevPoincareBall}).  Write
\[
q=\frac n{n-1},\qquad
\lambda=n+\frac n{n-1}\Bigl(1-\frac\tau p\Bigr).
\]
Then there is $N'\ge0$, depending only on $n$, $p$, $\tau$, $N$ and $d$, with the
following property.

Let $\Omega,\Omega',\Omega''\subseteq\R^n$ be open with $\Omega'$ bounded,
$\overline{\Omega'}\subseteq\Omega$ and $\overline{\Omega''}\subseteq\Omega'$,
and suppose
\[
d\le|x-y|\qquad\text{for every }x\in\Omega''\text{ and every }y\notin\Omega'
\]
(that is, $d$ is a positive lower bound for $d(\Omega'',\partial\Omega')$).  Let
$u\in W^{1,p}_{\mathrm{loc}}(\Omega)$ with weak gradient $g=\nabla u$
(\noderef{MemW1pLocOn}), and suppose
\[
\|\,|\nabla u|\,\|_{L^{p,n-\tau}(\Omega')}\le N
\qquad\text{and}\qquad
\|u\|_{L^{q}(\Omega')}\le N
\]
(\noderef{MorreyBoundOn}; the second is the case $\lambda=0$ of that predicate,
which is the bound $\int_{\Omega'(x,\rho)}|u|^q\,dx\le N^q$ for all
$x\in\Omega'$ and $0<\rho<\operatorname{diam}\Omega'$, together with
$u\in L^q(\Omega')$).  Then
\[
\|u\|_{\mathcal L^{q,\lambda}(\Omega'')}\le N'
\]
(\noderef{CampanatoBoundOn}), i.e.\
$u\in\mathcal L^{\frac n{n-1},\,n+(\frac n{n-1})(1-\frac\tau p)}(\Omega'')$.

Here, as in \S4 of the paper, $\Omega'(x,\rho)=\Omega'\cap B_\rho(x)$ and
$\Omega''(x,\rho)=\Omega''\cap B_\rho(x)$, the Morrey bound on a set is the
bound on $\int_{\cdot(x,\rho)}|\cdot|^q$ over all centres in that set and all
$0<\rho<\operatorname{diam}$ of it, and the Campanato bound subtracts the mean
over $\Omega''(x_0,\rho)$.

This is Lemma~Lemma~4.2 of the thesis of the paper, whose Morrey and Campanato spaces are
the local ones (``Campanato and Morrey spaces have obvious local versions'').
The formal statement realises ``local'' on both sides, as the paper's statement
does: the hypotheses are Morrey bounds on $\Omega'$ and the conclusion is a
Campanato bound on an inner subdomain $\Omega''$ compactly contained in
$\Omega'$.  The paper's Lemma~Lemma~4.2 of the thesis asserts membership only --- it names
no constant and gives no dependency list.  The quantitative form above is the
form the paper's own route delivers (``an immediate consequence of the
Sobolev--Poincar\'e inequality'': that inequality carries a constant $C(n)$, and
tracking it through H\"older's inequality and the Morrey hypothesis produces a
bound of exactly this shape), and it is the form Theorems~Theorem~3 of the thesis(ii) and
Theorem~5 of the thesis(ii) need, since those state bounds on
$\|u\|_{\mathcal L^{q,\lambda}(\Omega')}$ rather than mere membership.  Within
that quantitative form, $N'$ is allowed to depend on the separation $d$ of
$\Omega''$ from $\partial\Omega'$ in addition to $n,p,\tau,N$.  That dependence is necessary and not an artefact: for
$\rho$ comparable to $\operatorname{diam}\Omega''$ the ball $B_\rho(x_0)$ need
not lie in $\Omega'$, the Sobolev--Poincar\'e inequality is then unavailable, and
the only surviving estimate is the one coming from $\|u\|_{L^q(\Omega')}\le N$,
which gives $\int_{\Omega''(x_0,\rho)}|u-u_{x_0,\rho}|^q\le 2^{q+1}N^q$ and
therefore requires $(N')^q\ge2^{q+1}N^q\rho^{-\lambda}$; since $\lambda>0$
throughout the declared parameter range, $\rho^{-\lambda}$ is unbounded as
$\rho\downarrow0$, so no bound uniform in the position of $\Omega''$ inside
$\Omega'$ is available.  Letting $N'$ depend on $d$ still yields a finite
Campanato norm on every subdomain compactly contained in $\Omega'$, which is
exactly what local membership asserts, so the statement remains at least as
strong as the paper's.
\end{lemma}
\begin{proof}
Throughout, $q=\frac n{n-1}$ and $\lambda=n+q(1-\frac\tau p)$.  Note first the
two elementary facts used below.

\emph{(0a) The exponent identity.}  $\lambda=q\bigl(n-\frac\tau p\bigr)$.
Indeed
\[
q\Bigl(n-\frac\tau p\Bigr)
=q\Bigl((n-1)+1-\frac\tau p\Bigr)
=q(n-1)+q\Bigl(1-\frac\tau p\Bigr)
=n+q\Bigl(1-\frac\tau p\Bigr)=\lambda ,
\]
since $q(n-1)=\frac n{n-1}(n-1)=n$.

\emph{(0b) $\lambda>0$.}  As $n\ge2$ we have $q=\frac n{n-1}>0$.  Also
$\tau>0$ and $p>1$ give $\frac\tau p<\tau\le n$, so $n-\frac\tau p>0$; by (0a),
$\lambda=q(n-\frac\tau p)>0$.

Let $C=C(n)>0$ be the constant supplied by
\noderef{CitedSobolevPoincareBall}, so that for every ball $B_\sigma(y)$ and
every $w\in W^{1,1}(B_\sigma(y))$ with weak gradient $g_w$,
\[
\Bigl(\int_{B_\sigma(y)}\bigl|w-w_{B_\sigma(y)}\bigr|^{q}dx\Bigr)^{1/q}
\le C\int_{B_\sigma(y)}|g_w|\,dx ,
\qquad
w_{B_\sigma(y)}=\frac1{|B_\sigma(y)|}\int_{B_\sigma(y)}w\,dx
\tag{$*$}
\]
(the exponent $\frac{n-1}n$ there is $1/q$).  Let $\omega_n=|B_1(0)|$, so
$|B_\sigma(y)|=\omega_n\sigma^{\,n}$.  Define
\[
K=C\,\omega_n^{1-\frac1p},
\qquad
N'=2^{\frac{q+1}q}\,N\,\max\bigl\{K,\;d^{-\lambda/q}\bigr\} .
\]
Then $N'\ge0$, and $N'$ depends only on $n,p,\tau,N,d$: indeed $C$ and
$\omega_n$ depend only on $n$, and $q$ and $\lambda$ only on $n,p,\tau$.  Fix
$\Omega,\Omega',\Omega'',u,g$ as in the statement.

\emph{Step 1: the oscillation lemma.}  We first record a fact used in both
regimes below.  Let $A\subseteq B$ be measurable sets with $0<|A|<\infty$, let
$w\in L^1(A)$ --- so that the mean $w_A=|A|^{-1}\int_Aw\,dx$ is defined --- and
let $c\in\R$ be arbitrary.  Then
\[
\int_A|w-w_A|^{q}\,dx\;\le\;2^{\,q+1}\int_B|w-c|^{q}\,dx .
\tag{1}
\]
Both sides are integrals of nonnegative measurable functions, so each lies in
$[0,\infty]$ and the inequality is to be read there; in particular it is trivial
if the right-hand side is $+\infty$, so we may assume
$\int_B|w-c|^{q}dx<\infty$, and then also $\int_A|w-c|^{q}dx<\infty$ since
$A\subseteq B$.  To see the inequality, note $c-w_A=|A|^{-1}\int_A(c-w)\,dx$,
hence
\[
|c-w_A|\le\frac1{|A|}\int_A|w-c|\,dx=:m .
\]
Therefore $|w(x)-w_A|\le|w(x)-c|+|c-w_A|\le|w(x)-c|+m$ for every $x\in A$, and
since $(s+t)^q\le2^q(s^q+t^q)$ for all $s,t\ge0$ and $q\ge1$ (here $q>1$),
\[
|w(x)-w_A|^{q}\le2^{q}\bigl(|w(x)-c|^{q}+m^{q}\bigr).
\]
Integrating over $A$,
\[
\int_A|w-w_A|^{q}dx\le2^{q}\Bigl(\int_A|w-c|^{q}dx+|A|\,m^{q}\Bigr).
\]
Here $m<\infty$ because $w\in L^1(A)$ and $|A|<\infty$.  By Jensen's inequality
for the probability measure $|A|^{-1}\,dx$ on $A$ and the convex function
$t\mapsto t^{q}$ on $[0,\infty)$ (convex because $q>1$),
\[
m^{q}=\Bigl(\frac1{|A|}\int_A|w-c|\,dx\Bigr)^{q}
\le\frac1{|A|}\int_A|w-c|^{q}dx ,
\qquad\text{so}\qquad
|A|\,m^{q}\le\int_A|w-c|^{q}dx .
\]
Hence $\int_A|w-w_A|^{q}dx\le2^{q}\cdot2\int_A|w-c|^{q}dx=2^{q+1}\int_A|w-c|^{q}dx$,
and since the integrand is nonnegative and $A\subseteq B$ this is at most
$2^{q+1}\int_B|w-c|^{q}dx$, which is (1).

\emph{Step 2: $u\in L^{q}(\Omega'')$.}  The hypothesis
$\|u\|_{L^{q}(\Omega')}\le N$ (\noderef{MorreyBoundOn} with exponent $0$)
includes $u\in L^{q}(\Omega')$, and $\Omega''\subseteq\overline{\Omega''}
\subseteq\Omega'$, so $u\in L^{q}(\Omega'')$.  This is the first half of the
assertion $\|u\|_{\mathcal L^{q,\lambda}(\Omega'')}\le N'$
(\noderef{CampanatoBoundOn}).

For the second half, fix $x_0\in\Omega''$ and $\rho$ with
$0<\rho<\operatorname{diam}\Omega''$, and write $A=\Omega''(x_0,\rho)=
\Omega''\cap B_\rho(x_0)$ and $u_A=|A|^{-1}\int_Au\,dx$.  We must show
\[
\int_A|u-u_A|^{q}dx\le (N')^{q}\rho^{\lambda}.
\tag{2}
\]
Two preliminary remarks.  First, $A$ is a nonempty open set --- it contains
$x_0$, and $\Omega''$ and $B_\rho(x_0)$ are open --- so $|A|>0$; and
$A\subseteq\Omega'$ with $\Omega'$ bounded gives $|A|<\infty$.  So (1) applies to
$A$.  Second, $\Omega''\subseteq\overline{\Omega''}\subseteq\Omega'$ and $\Omega'$
is bounded, so $\Omega''$ is bounded as well and both diameters are finite;
monotonicity of the diameter then gives
$\operatorname{diam}\Omega''\le\operatorname{diam}\Omega'$, hence
\[
0<\rho<\operatorname{diam}\Omega' ,
\tag{3}
\]
so both Morrey hypotheses are available at the centre $x_0$ (which lies in
$\Omega'$) and the radius $\rho$.

\emph{Step 3: the regime $\rho<d$.}  We claim $B_\rho(x_0)\subseteq\Omega'$.
Indeed, let $z\in B_\rho(x_0)$ and suppose $z\notin\Omega'$; the separation
hypothesis applied to $x=x_0\in\Omega''$ and $y=z$ gives $d\le|x_0-z|<\rho<d$, a
contradiction.  Consequently
\[
\Omega'(x_0,\rho)=\Omega'\cap B_\rho(x_0)=B_\rho(x_0),
\tag{4}
\]
so the Morrey hypothesis on the gradient becomes a bound over the full ball.
Explicitly, that hypothesis is the Morrey bound of \noderef{MorreyBoundOn} for
the scalar function $x\mapsto|\nabla u(x)|$ with exponents $p$ and $n-\tau$ on
$\Omega'$ (note $n-\tau\ge0$ since $\tau\le n$, and
$\bigl||\nabla u|\bigr|^p=|\nabla u|^p$), so by (3) and $x_0\in\Omega'$ it reads
$\int_{\Omega'(x_0,\rho)}|\nabla u|^p\,dx\le N^p\rho^{\,n-\tau}$; with (4) this
is
\[
\int_{B_\rho(x_0)}|\nabla u|^{p}dx\le N^{p}\rho^{\,n-\tau}.
\tag{5}
\]

Next, $u\in W^{1,1}(B_\rho(x_0))$ with weak gradient $g$.  To see this, note
$\overline{B_\rho(x_0)}\subseteq\overline{\Omega'}\subseteq\Omega$ by (4), and
$\overline{B_\rho(x_0)}$ is compact, so the local integrability clause of
$u\in W^{1,p}_{\mathrm{loc}}(\Omega)$ (\noderef{MemW1pLocOn}) gives
$u\in L^{p}(\overline{B_\rho(x_0)})$ and $|g|\in L^{p}(\overline{B_\rho(x_0)})$;
as $|B_\rho(x_0)|<\infty$ and $p>1$, H\"older's inequality upgrades these to
$u\in L^{1}(B_\rho(x_0))$ and $|g|\in L^{1}(B_\rho(x_0))$.  The weak-gradient
identity of \noderef{MemW1pLocOn} holds for every test function smooth with
compact support inside $\Omega$, in particular for every test function with
compact support inside $B_\rho(x_0)\subseteq\Omega$, and for such a test function
the integrals over $\Omega$ and over $B_\rho(x_0)$ agree because the test
function vanishes off $B_\rho(x_0)$; so $g$ is the weak gradient of $u$ on
$B_\rho(x_0)$ as well.

Now apply $(*)$ on $B_\rho(x_0)$ to $w=u$:
\[
\Bigl(\int_{B_\rho(x_0)}\bigl|u-u_{B_\rho(x_0)}\bigr|^{q}dx\Bigr)^{1/q}
\le C\int_{B_\rho(x_0)}|\nabla u|\,dx .
\]
Estimate the right-hand side by H\"older's inequality with exponents $p$ and
$\frac p{p-1}$ and then by (5):
\[
\int_{B_\rho(x_0)}|\nabla u|\,dx
\le|B_\rho(x_0)|^{1-\frac1p}\Bigl(\int_{B_\rho(x_0)}|\nabla u|^{p}dx\Bigr)^{1/p}
\le\bigl(\omega_n\rho^{\,n}\bigr)^{1-\frac1p}\bigl(N^{p}\rho^{\,n-\tau}\bigr)^{1/p} .
\]
The right-hand side equals
$\omega_n^{1-\frac1p}N\,\rho^{\,n(1-\frac1p)+\frac{n-\tau}p}
=\omega_n^{1-\frac1p}N\,\rho^{\,n-\frac\tau p}$, because
$n-\frac np+\frac np-\frac\tau p=n-\frac\tau p$.  Combining the last two
displays and raising to the power $q>0$,
\[
\int_{B_\rho(x_0)}\bigl|u-u_{B_\rho(x_0)}\bigr|^{q}dx
\le\bigl(C\,\omega_n^{1-\frac1p}N\bigr)^{q}\rho^{\,q(n-\frac\tau p)}
=(KN)^{q}\rho^{\lambda},
\]
the last step by (0a) and the definition of $K$.  Finally apply (1) with
$A=\Omega''(x_0,\rho)$, $B=B_\rho(x_0)$ and the constant
$c=u_{B_\rho(x_0)}$.  This is legitimate: $A=\Omega''\cap B_\rho(x_0)\subseteq
B_\rho(x_0)=B$, we have $0<|A|<\infty$ as noted above, and
$u\in L^{1}(B_\rho(x_0))$ was established in the previous paragraph, hence
$u\in L^1(A)$.  We obtain
\[
\int_A|u-u_A|^{q}dx
\le2^{\,q+1}\int_{B_\rho(x_0)}\bigl|u-u_{B_\rho(x_0)}\bigr|^{q}dx
\le2^{\,q+1}(KN)^{q}\rho^{\lambda}
=\bigl(2^{\frac{q+1}q}KN\bigr)^{q}\rho^{\lambda}
\le(N')^{q}\rho^{\lambda},
\]
since $2^{\frac{q+1}q}KN\le N'$ by the definition of $N'$.  This is (2) in this
regime.

\emph{Step 4: the regime $\rho\ge d$.}  Here $A\subseteq\Omega'$ and
$u\in L^{q}(\Omega')$ by Step 2, so $u\in L^{q}(A)$, and since $|A|<\infty$ and
$q>1$ H\"older's inequality gives $u\in L^{1}(A)$.  Apply (1) with
$A=\Omega''(x_0,\rho)$, $B=A$ and $c=0$:
\[
\int_A|u-u_A|^{q}dx\le2^{\,q+1}\int_A|u|^{q}dx .
\]
Since $A\subseteq\Omega'\cap B_\rho(x_0)=\Omega'(x_0,\rho)$ and the integrand is
nonnegative, and since by (3) the Morrey hypothesis
$\|u\|_{L^{q}(\Omega')}\le N$ (\noderef{MorreyBoundOn} with exponent $0$) applies
at the centre $x_0\in\Omega'$ and radius $\rho$,
\[
\int_A|u|^{q}dx\le\int_{\Omega'(x_0,\rho)}|u|^{q}dx\le N^{q}\rho^{0}=N^{q}.
\]
Hence $\int_A|u-u_A|^{q}dx\le2^{\,q+1}N^{q}$.  On the other hand $\rho\ge d>0$
and $\lambda>0$ by (0b), so $\rho^{\lambda}\ge d^{\lambda}$ and therefore
\[
\bigl(2^{\frac{q+1}q}N\,d^{-\lambda/q}\bigr)^{q}\rho^{\lambda}
=2^{\,q+1}N^{q}d^{-\lambda}\rho^{\lambda}
\ge2^{\,q+1}N^{q}.
\]
Since $2^{\frac{q+1}q}N\,d^{-\lambda/q}\le N'$ by the definition of $N'$, we get
$\int_A|u-u_A|^{q}dx\le2^{\,q+1}N^{q}\le(N')^{q}\rho^{\lambda}$, which is (2) in
this regime.

Steps 3 and 4 cover all $\rho$ with $0<\rho<\operatorname{diam}\Omega''$, so (2)
holds for every such $\rho$ and every $x_0\in\Omega''$.  Together with Step 2
this is precisely $\|u\|_{\mathcal L^{q,\lambda}(\Omega'')}\le N'$
(\noderef{CampanatoBoundOn}), with $N'$ depending only on $n,p,\tau,N$ and $d$.
\end{proof}
